% pub-dialogue-analysis.tex
%
% Analysis write-up: Methods and Results (v19 canonical run)
% Based on 2026-05-dialogue-results-sketch.tex
%
% NOTEBOOK CROSS-REFERENCES
% =========================
% Each Methods and Results section below has a comment indicating which
% notebook(s) implement or visualise that section.
%
% Notebook index:
%   00_data_quality.ipynb    — Assess-phase quality checks (question-agnostic)
%   01_processing.ipynb      — Full pipeline: ingest → chunk → extract → embed
%   01a_clustering.ipynb     — KMeans clustering, labelling, framing lenses
%   02_shared_structure.ipynb — R1: shared cross-cutting grammar
%   03_ai_distinctiveness.ipynb — R2: AI vs non-AI distinctiveness
%   04_temporal_dynamics.ipynb  — R3: temporal dynamics of AI concerns
%   05_robustness.ipynb         — M6: validation, k-sensitivity, prompt sensitivity
%

\textbf{Paper sketch: Methods and Results (v19)}

All numbers in this document come from the v19 run outputs
in~public\_dialogue\_analysis\_outputs-8.zip~(n=7,671 chunks). Where a
previous draft is referenced, it is for comparison only --- the analytic
claims rest on v19.

\textbf{Methods}

% ─────────────────────────────────────────────────────────────────────────────
% M1. CORPUS
% Notebook: 01_processing.ipynb — §Purpose and §Corpus (opening markdown cell)
%           00_data_quality.ipynb — corpus-level counts and technology breakdown
% ─────────────────────────────────────────────────────────────────────────────
\textbf{M1. Corpus}

\begin{itemize}
\item
  \textbf{66 UK public dialogue reports}~on emerging
  technologies,~\textbf{2004--2025}.
\item
  \textbf{11 technology categories}, with metadata mapping each document
  to a primary technology and publication year.
\item
  \textbf{Heavy AI weighting}: 41 of 66 documents are AI dialogues,
  reflecting recent commissioning patterns. Smaller corpora for genetic
  technologies (7 documents), gene editing (4), nanotechnology (4),
  nuclear (3), GM (2). One document each for carbon capture, drones,
  geoengineering, neuraltech, quantum.
\item
  \textbf{Year coverage is uneven}: substantial document gaps in
  2005--2006, 2011--2014, and 2020 (the latter driven by COVID-related
  interruption of public-dialogue commissioning). These gaps are
  reported transparently in the temporal analysis.
\item
  Corpus is imbalanced; smaller technology corpora have limited
  representation. Implications for analysis are discussed in
  Limitations.
\end{itemize}

\textbf{Table for Methods}: documents and chunks per technology, with
year ranges, drawn
from~paragraph\_chunks.csv~and~paragraph\_chunks\_per\_document.csv.

% ─────────────────────────────────────────────────────────────────────────────
% M2. SEGMENTATION
% Notebook: 01_processing.ipynb — §SECTION 1: Data ingestion and preprocessing
%           (cells 22–24: hybrid chunking, three-case strategy, chunk statistics)
%           00_data_quality.ipynb — contamination check, chunk-quality diagnostic
% ─────────────────────────────────────────────────────────────────────────────
\textbf{M2. Segmentation}

PDF text extracted via PyMuPDF. The chunker uses a three-case hybrid
strategy:

\begin{itemize}
\item
  \textbf{Case 1 --- paragraph segmentation alone (n=8 documents)}:
  paragraph-level splitting (on double newlines) produces ≥3 substantive
  paragraphs AND no paragraph exceeds the 500-word cap. All chunks are
  paragraphs as the author wrote them.
\item
  \textbf{Case 2 --- paragraphs with internal sentence-splitting (n=43
  documents)}: paragraph splitting produces enough substantive
  paragraphs but at least one would exceed the 500-word cap. Well-sized
  paragraphs are kept as-is; only the oversized paragraphs are split at
  sentence boundaries into sub-chunks of approximately 300 words. The
  author's paragraph boundaries are preserved everywhere they exist. (Of
  these 43, four documents are an extreme case in
  which~\emph{every}~paragraph is oversized --- primarily long-form
  survey reports.)
\item
  \textbf{Case 3 --- full sentence-fallback (n=15 documents)}: paragraph
  splitting produces fewer than 3 substantive paragraphs, indicating the
  PDF lacks paragraph break encoding that PyMuPDF can extract. The whole
  document is sentence-split and repacked into \textasciitilde300-word
  chunks. These are predominantly recent UK AI policy reports produced
  from HTML originals (e.g., Ada Lovelace Institute briefings, ONS
  surveys, BritainThinks/CDEI tracker reports), plus the 2004 Royal
  Society Nanotechnology report.
\end{itemize}

Each chunk records which method produced it via
the~chunking\_method~column (paragraph,~sentence\_split,
or~sentence\_fallback).

\textbf{Filters}: chunks below 40 words or 80 characters are dropped.
Chunks above 500 words are split at sentence boundaries before reaching
this filter, so only 5 chunks (0.07\% of corpus) have hit the truncation
cap as a residual safety net.

\textbf{Final corpus}: 7,671 chunks across 66 documents. Of
these,~\textbf{5,285 (69\%) are paragraph segments}, 1,675 (22\%) are
sentence-split sub-chunks from oversized paragraphs in otherwise
well-chunked documents, and 711 (9\%) come from full sentence-fallback.
Mean chunk length is 173 words; median is 138 words.

\textbf{Diagnostic note}: 832 of 7,671 chunks (10.8\%) are flagged by an
automated diagnostic as likely-bibliography (134) or likely-survey-table
(723) fragments. LLM extraction generally returns NO\_CONCERN for these,
limiting their distortive effect on findings.

% ─────────────────────────────────────────────────────────────────────────────
% M3. PHRASE EXTRACTION
% Notebook: 01_processing.ipynb — §Concern extraction (cells 30–36)
%                                  §SECTION 3B: Benefits (cells 44–50)
%           The word-boundary regex post-filter is in pub_dialogue/address.py
% ─────────────────────────────────────────────────────────────────────────────
\textbf{M3. Phrase extraction}

\begin{itemize}
\item
  Two parallel pipelines:~\textbf{concerns}~and~\textbf{benefits}.
\item
  LLM (gpt-5-mini) prompted to extract 1--3 decontextualised concern (or
  benefit) phrases per chunk, removing technology-specific language to
  enable cross-technology comparison.
\item
  Word-boundary regex post-filter to catch phrases retaining
  technology-specific terms.
\item
  \textbf{Output}: 17,596 concern phrases from 6,323 chunks (82.4\%
  chunk yield); 18,346 benefit phrases from 6,910 chunks (90.1\% chunk
  yield).
\item
  \textbf{Coverage}: 65 of 66 documents produced concerns; 65 of 66
  produced benefits. The single document producing neither
  (2009\_FSA\_GM\_Food.pdf) is a 1-chunk fragment.
\item
  \textbf{AI-vs-non-AI split}: of 17,596 concern phrases, AI contributes
  6,955 (39.5\%) and non-AI technologies contribute 10,641 (60.5\%).
\end{itemize}

\textbf{Reflexive note}: extraction frames "concerns" and "benefits" as
analytic categories. The choice shapes what surfaces. Sensitivity to
this choice is tested in M6.

% ─────────────────────────────────────────────────────────────────────────────
% M4. EMBEDDING, CLUSTERING, FRAMING LENSES
% Notebook: 01a_clustering.ipynb — §SECTION 3/4/5 (concerns) and 3B/4B/5B (benefits)
%   - KMeans at k=75 implemented in §SECTION 3 (cells 19–25)
%   - Cluster labelling (LLM, without technology metadata) in §SECTION 4 (cells 31–33)
%   - Framing lens grouping in §SECTION 5 (cells 37–41)
%   - k-sensitivity (k=60 and k=90) in 05_robustness.ipynb §SECTION 10/10B
% ─────────────────────────────────────────────────────────────────────────────
\textbf{M4. Embedding, clustering, framing lenses}

\begin{itemize}
\item
  Phrases embedded with~text-embedding-3-large; L2-normalised.
\item
  k-means clustering at k=75 (sensitivity tested at k=60 and k=90).
\item
  Each cluster characterised by size, technology distribution, and
  normalised entropy as a proxy for cross-cutting-ness.
\item
  \textbf{Cross-cutting threshold}: normalised entropy ≥ 0.5.
\item
  LLM-assigned cluster labels (without access to source-technology
  metadata, to avoid labelling contamination).
\item
  Clusters grouped into~\textbf{10 framing lenses for
  concerns}~and~\textbf{10 framing lenses for benefits}~by a separate
  LLM step. Lenses are conceptual organisations of the cluster space;
  their empirical validity is tested via dendrogram analysis (see M6).
\end{itemize}

% ─────────────────────────────────────────────────────────────────────────────
% M5. COMPARISON METRICS
% Notebook: 03_ai_distinctiveness.ipynb
%   - Technology-weighted baseline: cells 2–8
%   - Document-weighted baseline (diagnostic): cell 4
%   - Lens-level and cluster-level difference calculations: cells 6–16
% ─────────────────────────────────────────────────────────────────────────────
\textbf{M5. Comparison metrics}

\begin{itemize}
\item
  \textbf{Two baseline constructions}~for AI-vs-non-AI comparison:

  \begin{itemize}
  \item
    \emph{Technology-weighted}: each non-AI technology contributes
    equally.
  \item
    \emph{Document-weighted}: non-AI phrases pooled, weighted by their
    natural document-volume distribution.
  \end{itemize}
\item
  Comparing both baselines is a robustness check: findings that hold
  under both are not artefacts of weighting choice.
\item
  \textbf{Distinctiveness measured at two levels}: cluster-level
  (per-cluster AI−non-AI share difference) and lens-level (per-lens
  difference).
\end{itemize}

% ─────────────────────────────────────────────────────────────────────────────
% M6. VALIDATION AND SENSITIVITY
% Notebook: 05_robustness.ipynb
%   (a) Human inter-rater validation: cell 4 (stratified sample generation)
%   (b) Prompt sensitivity: §SECTION 10C (cell 31) — V0/V1/V3 variants
%   (c) Lens-scheme dendrogram: 02_shared_structure.ipynb cells 4/18
%                               (ARI computed there; reported here)
%   (d) k-sensitivity: 05_robustness.ipynb §SECTION 10/10B (cells 16–20, 30)
% ─────────────────────────────────────────────────────────────────────────────
\textbf{M6. Validation and sensitivity}

Four validation strands:

\textbf{(a) Human inter-rater validation}~on a stratified random sample
of 50 paragraphs (expanded in final canonical run pending). Cohen's κ on
the binary "concern present" decision = 0.84 (almost-perfect agreement
on the Landis--Koch scale). 92\% raw agreement.~\textbf{Asymmetric
error}: LLM prone to over-extraction (4 cases where the LLM extracted a
concern but the human coder did not) and not under-extraction (0 cases
where the human coder identified a concern the LLM missed). Procedural
and methodological text occasionally mined for concerns it doesn't
substantively contain.

\textbf{(b) Prompt sensitivity}~on a stratified sample of 529 paragraphs
under three prompt variants: V0 (baseline), V1 (broader "issues"
framing), V3 (strict decontextualisation). The qualitative pattern of AI
distinctiveness was preserved across all three variants. The Privacy
effect's magnitude was sensitive to decontextualisation strictness (V0:
+21pp; V1: +16pp; V3: +9pp), suggesting V0 is an upper bound and V3 a
lower bound on the precise pp difference.~\emph{(Note: the prompt
sensitivity analysis was run against the v16 corpus; rerunning against
v19 chunks is a candidate task before submission, though qualitative
direction is expected to hold.)}

\textbf{(c) Lens-scheme dendrogram validation}: hierarchical clustering
of cluster centroids compared against LLM-derived framing lens scheme
via Adjusted Rand Index. Concerns ARI = 0.185; benefits ARI = 0.150. The
lens scheme is an interpretively useful organisation, not a strong
empirical partition. Some lenses (notably Privacy/Consent \& Data)
validate cleanly; others are looser conceptual groupings. Reported
findings are hedged accordingly.

\textbf{(d) k-sensitivity}~at k ∈ \{60, 75, 90\}. Headline findings
preserved across all three.

\textbf{Supplementary Table S1}: summary of validation, prompt
sensitivity, and dendrogram results.

% ─────────────────────────────────────────────────────────────────────────────
% M7. NOTES ON DATA QUALITY
% Notebook: 00_data_quality.ipynb — full corpus quality report
%   - Contamination rate (bibliography/table flags): cells 11–15
%   - Truncation elimination in v19: reported in 01_processing.ipynb cell 22–24
%   - 2020 gap: visible in temporal plots in 04_temporal_dynamics.ipynb
% ─────────────────────────────────────────────────────────────────────────────
\textbf{M7. Notes on data quality}

\begin{itemize}
\item
  \textbf{Truncation is effectively eliminated}~in v19 (5 chunks,
  0.07\%, residual). Long paragraphs are sentence-split rather than
  truncated; documents lacking paragraph breaks go through full
  sentence-fallback. This is a substantive change from earlier pipeline
  iterations.
\item
  \textbf{The 2020 year contains only 31 chunks}~(one fragmentary
  document). COVID-related interruption of public-dialogue commissioning
  is the likely cause. Worth noting in the temporal analysis.
\item
  \textbf{Chunk-quality contamination is 10.8\%}~(8\% of
  likely-table-rows, 2\% likely-bibliography). The automated diagnostic
  doesn't filter them; the LLM extraction step generally returns
  NO\_CONCERN, so distortive effect on findings is limited.
\end{itemize}

\textbf{Results}

\textbf{Headlines}~(all from v19 outputs):

UK public dialogue reports repeatedly frame emerging technologies
through procedural legitimacy: governance, accountability, transparency,
public participation, trust, evidence, oversight, and distribution of
benefits/harms. AI does not sit outside this shared structure but
expresses it through informational concerns: misuse of personal data,
data transparency, mistrust of institutional data handling, privacy and
security, control over personal data, and loss of human oversight. AI
benefits are framed less around environmental and physical-world
transformation and more around service efficiency, data privacy and
control, research quality, and personalisation. AI concern profiles stay
broadly distributed over time, while the content mix shifts across
commissioning periods.

% ─────────────────────────────────────────────────────────────────────────────
% R1. SHARED CROSS-CUTTING GRAMMAR OF LEGITIMACY
% Notebook: 02_shared_structure.ipynb
%   - Normalised entropy and cross-cutting threshold: cells 6–10
%   - Stable-core scatterplot (Figure 1): cell 12
%   - Top-20 cross-cutting clusters table (Table 1): cells 8–10
%   - Benefit parallel (67% cross-cutting): cells 20–26
%   - PCA embedding space visualisation: cells 14, 28
% ─────────────────────────────────────────────────────────────────────────────
\textbf{R1. Public concerns about emerging technology share a procedural
grammar of legitimacy}

\textbf{What it claims}: across the corpus, the dominant pattern is
shared concern about how technology is governed, who participates in
decisions, and how information flows. The substance of the technology
matters less than the procedural questions surrounding it.

\textbf{Empirical substantiation}:

\begin{itemize}
\item
  Of 75 concern clusters,~\textbf{62 (83\%) are
  cross-cutting}~(normalised entropy ≥ 0.5).
\item
  \textbf{Cross-cutting clusters hold 14,771 of 17,596 concern phrases
  (83.9\%)}.
\item
  The 20 largest cross-cutting clusters together account for 40.1\% of
  all concern phrases. They are about procedural and structural
  concerns: misuse of personal data, distrust in governance and
  oversight, lack of transparency in decision-making, lack of oversight
  and accountability, uncertainty about future impacts, inadequate
  testing and reliability, ethical and moral boundaries, insufficient
  public engagement, safety risks, insufficient public participation in
  decisions, unclear and inconsistent regulation, limited public
  influence on research, short-term fragmented responses, erosion of
  social connection, insufficient human oversight, insufficient support
  and infrastructure.
\item
  The benefit side mirrors this:~\textbf{50 of 75 benefit clusters
  (67\%) are cross-cutting}, holding 13,190 of 18,346 benefit phrases
  (71.9\%). The benefit side has a higher proportion of tech-specific
  clusters because medical and clinical benefits cluster around genetic
  technologies, and service-efficiency benefits cluster around AI.
\end{itemize}

\textbf{Figure 1}: stable-core scatterplot of the 75 concern clusters
(normalised entropy × frequency), with the cross-cutting +
high-frequency quadrant labelled as the "shared core" and annotated with
the top cluster labels.

\includegraphics[width=6.26806in,height=4.41181in]{media/image1.png}

\textbf{Table 1}: top 20 cross-cutting concern clusters by global
prevalence.

Top 20 cross-cutting concern clusters (normalised entropy ≥ 0.5), ranked
by share of all concern phrases. These 20 clusters account for 40.1\% of
all 17,596 extracted concern phrases across 65 documents and 11
technologies.

\begin{longtable}[]{@{}lllll@{}}
\toprule
\textbf{Cluster ID} & \textbf{Cluster label} & \textbf{N phrases} &
\textbf{Share of all concerns} & \textbf{Tech entropy (raw)} \\
\midrule
\endhead
50 & Misuse of personal data & 531 & 3.02\% & 0.78 \\
62 & Distrust in governance and oversight & 512 & 2.91\% & 1.83 \\
46 & Lack of transparency in decision-making & 502 & 2.85\% & 1.67 \\
12 & Lack of oversight and accountability & 390 & 2.22\% & 1.55 \\
43 & Uncertainty about future impacts & 375 & 2.13\% & 1.90 \\
64 & Privacy and data security & 361 & 2.05\% & 0.91 \\
65 & Inadequate testing and reliability & 358 & 2.03\% & 1.50 \\
0 & Ethical and moral boundaries & 354 & 2.01\% & 1.16 \\
4 & Insufficient public engagement & 328 & 1.86\% & 1.87 \\
67 & Safety, security and environmental risks & 326 & 1.85\% & 1.95 \\
3 & Insufficient public participation in decisions & 322 & 1.83\% &
1.69 \\
36 & Unclear and inconsistent regulation & 316 & 1.80\% & 1.79 \\
25 & Inaccurate diagnosis and treatment & 314 & 1.78\% & 1.34 \\
39 & Limited public influence on research & 312 & 1.77\% & 1.37 \\
73 & Short-term, fragmented responses & 311 & 1.77\% & 1.95 \\
14 & Erosion of social connection and identity & 301 & 1.71\% & 1.53 \\
63 & Insufficient human oversight & 295 & 1.68\% & 1.24 \\
52 & Insufficient support and infrastructure & 285 & 1.62\% & 1.71 \\
16 & Unclear purpose and decision-making & 284 & 1.61\% & 1.73 \\
54 & Unknown and unpredictable risks & 283 & 1.61\% & 1.69 \\
\bottomrule
\end{longtable}

\textbf{Reading the finding}: this is a baseline finding, not a
"discovery". UK public dialogues commissioned across two decades and 11
technologies converge on shared procedural concerns. The 20 most
prominent shared concerns are dominated by governance, oversight,
accountability, transparency, public engagement, and harm-distribution
themes. Two themes specifically about data (Misuse of personal data and
Privacy and data security) are also in the top 20 --- this is the first
signal of AI's informational footprint, but it sits~\emph{within}~a
shared procedural grammar that runs across all 11 technologies. AI's
distinctiveness, when it appears, expresses this grammar through
informational concerns rather than departing from it.

% ─────────────────────────────────────────────────────────────────────────────
% R2. AI DISTINCTIVENESS: INFORMATIONAL AUTONOMY
% Notebook: 03_ai_distinctiveness.ipynb
%   - Technology-weighted lens comparison (radar): cells 2–8, 12
%   - Document-weighted baseline diagnostic: cell 4
%   - Cluster-level distinctiveness bar chart (Figure 3): cell 38
%   - AI vs genetic technologies tech-specific clusters: cells 14–16
%   - Benefit mirror (radar, Table 2 panel B): cells 20–30
%   - Lens-level distinctiveness table (Table 2 panel B): cells 28–30
% ─────────────────────────────────────────────────────────────────────────────
\textbf{R2. AI public dialogue is distinctively focused on data and
informational autonomy}

\textbf{What it claims}: within the shared grammar, AI public dialogue
carries a distinctive informational footprint. AI concerns are
disproportionately about data, privacy, transparency, control over
personal information, and the displacement of human oversight and
contact. AI is correspondingly under-represented in concerns about
physical safety, environmental harm, and embodied risk.

\textbf{Empirical substantiation, three layers of evidence}:

\textbf{(i) Lens-level distinctiveness}~(the headline numbers). AI
is~\textbf{+16.4pp on the Privacy, Consent \& Data lens}~(21.9\% of AI
concerns vs 5.6\% of non-AI under the technology-weighted baseline;
identical direction at +16.7pp under the document-weighted baseline). AI
is over-indexed on Security, Misuse \& Threats (+4.6pp tech-weighted),
Equity \& Inclusion (+3.2pp), System Readiness \& Reliability (+2.6pp),
and Transparency \& Communication (+2.3pp). AI is
correspondingly~\textbf{under-indexed on Safety, Health \& Environment
(-13.3pp tech-weighted, -15.3pp doc-weighted)}~and on Purpose, Benefits
\& Economic Impact (-3.0pp), Public Participation \& Engagement
(-2.7pp), and Governance \& Accountability (-2.6pp).~\textbf{Both
baselines tell the same rank-order story; magnitudes differ slightly.}

\textbf{(ii) Cluster-level distinctiveness}: the top 10 most
AI-over-indexed concern clusters are dominated by data and informational
themes.~\textbf{Seven of the top 10}~are explicitly about data or human
oversight: Misuse of personal data (+5.0pp), Lack of data transparency
(+3.0pp), Mistrust of government data handling (+2.8pp), Privacy and
data security (+2.6pp), Loss of human oversight and contact (+2.1pp),
Insufficient human oversight (+2.0pp), Loss of control over personal
data (+1.9pp). The remaining three concern equity and accountability
themes that the AI corpus emphasises (Discrimination and unequal
impacts, Lack of oversight and accountability, Unequal public
awareness). The least AI-distinctive clusters are about physical and
environmental harm: Safety/security/environmental risks (-2.1pp),
Environmental and ecosystem impacts (-2.1pp), Environmental
contamination (-1.8pp), Unclear and inconsistent regulation (-1.5pp),
Uncertain environmental and community impacts (-1.4pp).

\textbf{(iii) Contrast with genetic technologies}. AI is one of two
technologies in the corpus producing distinctive technology-specific
cluster footprints; the other is genetic technologies. The two
footprints occupy non-overlapping domains. AI's~\textbf{eight
tech-specific clusters}~are all about informational autonomy or labour:
Misuse of personal data, Privacy and data security, Lack of data
transparency, Mistrust of government data handling, Loss of control over
personal data, Loss of human oversight and contact, Unequal public
awareness, and two on job displacement. Genetic
technologies'~\textbf{three tech-specific clusters}~are all about
embodied or research ethics: Ethical and moral boundaries, Emotional
burden on families, and Ethical boundaries in research. Nuclear produces
one tech-specific cluster, Long-term waste storage safety. No other
technology produces tech-specific clusters of meaningful size. This
contrast supports the interpretation that AI's distinctiveness
is~\emph{substantively}~informational rather than reflecting whatever
themes happen to surface in higher-volume technology corpora.

\textbf{The mirror-image finding on benefits}: AI is over-indexed on
Service Efficiency \& Personalisation (+13.9pp), Data Privacy \& Control
(+7.4pp, the protective mirror to the concern about data misuse),
Research Quality \& Ethics (+4.7pp), and Equity, Inclusion \& Wellbeing
(+2.5pp). AI is~\textbf{under-indexed on Safety, Risk \& Environment
(-15.1pp)}~and Economic \& Local Benefits (-5.7pp), Information \&
Communication (-4.3pp). The benefit fingerprint mirrors the concern
fingerprint: AI is associated with promised informational and
service-efficiency gains, much less with environmental or physical-world
benefits.

\textbf{Robustness}: the qualitative direction of these findings holds
across the three prompt variants tested (V0/V1/V3); the magnitude on
Privacy varies by approximately a factor of two between V0 and V3. The
headline interpretation --- AI public dialogue is distinctively about
informational autonomy --- is robust; the precise pp difference is not.
Both technology-weighted and document-weighted baselines produce the
same rank order of distinctive lenses for both concerns and benefits.

\textbf{Figure 2}: AI vs non-AI lens radar chart for concerns,
visualising the asymmetry around Privacy/Consent/Data and the
corresponding under-indexing on Safety/Health/Environment.

\includegraphics[width=6.26806in,height=3.33958in]{media/image2.emf}

\textbf{Figure 3}: AI vs non-AI cluster-level distinctiveness,
horizontal bar chart showing the 20 most distinctive concern clusters
(top 10 over-indexed and top 10 under-indexed) with AI share and non-AI
share visible side by side.

\includegraphics[width=6.26806in,height=5.07431in]{media/image3.png}

\textbf{Table 2}: top AI-distinctive concern clusters and framing
lenses.

\textbf{Panel A --- cluster-level distinctiveness}. Top 10 most
AI-over-indexed and top 10 most AI-under-indexed concern clusters,
ranked by percentage-point difference between AI share and
(technology-weighted) non-AI share.

\begin{longtable}[]{@{}lllll@{}}
\toprule
\textbf{Cluster label} & \textbf{AI share (\%)} & \textbf{Non-AI share
(\%)} & \textbf{Difference (pp)} & \textbf{Direction} \\
\midrule
\endhead
Misuse of personal data & 5.9 & 0.9 & +5.0 & AI over-indexed \\
Lack of data transparency & 3.2 & 0.2 & +3.0 & AI over-indexed \\
Mistrust of government data handling & 2.9 & 0.1 & +2.8 & AI
over-indexed \\
Privacy and data security & 3.9 & 1.3 & +2.6 & AI over-indexed \\
Loss of human oversight and contact & 2.1 & 0.0 & +2.1 & AI
over-indexed \\
Insufficient human oversight & 2.8 & 0.9 & +2.0 & AI over-indexed \\
Loss of control over personal data & 2.0 & 0.1 & +1.9 & AI
over-indexed \\
Discrimination and unequal impacts & 2.2 & 0.7 & +1.4 & AI
over-indexed \\
Lack of oversight and accountability & 2.8 & 1.6 & +1.2 & AI
over-indexed \\
Unequal public awareness & 1.5 & 0.4 & +1.1 & AI over-indexed \\
Safety, security and environmental risks & 1.4 & 3.5 & −2.1 & AI
under-indexed \\
Environmental and ecosystem impacts & 0.2 & 2.3 & −2.1 & AI
under-indexed \\
Environmental contamination and release & 0.1 & 1.9 & −1.8 & AI
under-indexed \\
Unclear and inconsistent regulation & 0.7 & 2.2 & −1.5 & AI
under-indexed \\
Uncertain environmental and community impacts & 0.2 & 1.6 & −1.4 & AI
under-indexed \\
High costs and affordability & 1.0 & 2.4 & −1.4 & AI under-indexed \\
Short-term, fragmented responses & 1.3 & 2.7 & −1.4 & AI
under-indexed \\
Misuse by malicious actors & 1.0 & 2.3 & −1.3 & AI under-indexed \\
Information complexity and comprehension & 1.5 & 2.8 & −1.2 & AI
under-indexed \\
Long-term health and environmental impacts & 0.1 & 1.4 & −1.2 & AI
under-indexed \\
\bottomrule
\end{longtable}

\textbf{Panel B --- lens-level distinctiveness (both baselines)}.
Per-lens AI-vs-non-AI share difference under technology-weighted and
document-weighted baselines.

\begin{longtable}[]{@{}llllll@{}}
\toprule
\textbf{Framing lens} & \textbf{AI share (\%)} & \textbf{Non-AI share,
tech-weighted (\%)} & \textbf{Δ tech-weighted (pp)} & \textbf{Non-AI
share, doc-weighted (\%)} & \textbf{Δ doc-weighted (pp)} \\
\midrule
\endhead
Privacy, Consent \& Data & 21.9 & 5.6 & +16.4 & 5.3 & +16.7 \\
Security, Misuse \& Threats & 18.0 & 13.4 & +4.6 & 9.3 & +8.7 \\
Equity \& Inclusion & 18.5 & 15.4 & +3.2 & 13.8 & +4.7 \\
System Readiness \& Reliability & 19.3 & 16.7 & +2.6 & 18.6 & +0.7 \\
Transparency \& Communication & 15.1 & 12.8 & +2.3 & 11.7 & +3.4 \\
Ethical \& Societal Boundaries & 8.0 & 7.9 & +0.1 & 11.5 & −3.4 \\
Governance \& Accountability & 24.3 & 26.9 & −2.6 & 28.1 & −3.7 \\
Public Participation \& Engagement & 13.8 & 16.5 & −2.7 & 16.7 & −2.9 \\
Purpose, Benefits \& Economic Impact & 10.6 & 13.5 & −3.0 & 12.5 &
−1.9 \\
Safety, Health \& Environment & 11.8 & 25.1 & −13.3 & 27.1 & −15.3 \\
\bottomrule
\end{longtable}

\textbf{Supplementary tables}: full lens-level distinctiveness under
both baselines for benefits; sensitivity numbers across V0/V1/V3
prompts.

% ─────────────────────────────────────────────────────────────────────────────
% R3. TEMPORAL DYNAMICS: STABLE SPREAD, SHIFTING CONTENT
% Notebook: 04_temporal_dynamics.ipynb
%   - Four-window entropy calculation: cells 4–8
%   - Content mix by time window (top-5 per window): cells 10–12
%   - Heatmap of cluster shares across windows (Figure 4): cell 14/28
%   - Benefit trajectory (rising/declining clusters): cells 22–26
%   - Year-by-year displacement in concern/benefit space: cells 6–8, 18–20
% ─────────────────────────────────────────────────────────────────────────────
\textbf{R3. AI public discourse content shifts over time within stable
spread}

\textbf{What it claims}: the spread of AI concerns across the cluster
space is stable across two decades; the content mix shifts in
interpretable ways that reflect both substantive evolution and the
changing landscape of dialogue commissioning.

\textbf{Empirical substantiation}:

\begin{itemize}
\item
  Across four time windows (2004--2017, 2018--2020, 2021--2023,
  2024--2025), the~\textbf{normalised entropy of AI concern
  distributions stays at \textasciitilde0.91--0.95}. The spread is
  stable across two decades.
\item
  Sample sizes per window: 1,432 phrases (2004--2017, n=5 documents);
  1,568 (2018--2020, n=6); 2,427 (2021--2023, n=19); 1,528 (2024--2025,
  n=11). The corpus is concentrated post-2017.
\item
  \textbf{Content mix shifts}:

  \begin{itemize}
  \item
    \textbf{2004--2017}: data themes dominate. Top 5 --- Misuse of
    personal data (8.6\%), Lack of data transparency (5.1\%), Privacy
    and data security (4.7\%), Loss of human oversight and contact
    (3.6\%), Invasion of personal privacy (3.2\%).
  \item
    \textbf{2018--2020}: data themes recede; governance, equity, and
    autonomy concerns rise. Top 5 --- Distrust in governance and
    oversight (3.4\%), Resource scarcity and inefficient use (3.2\%),
    Loss of personal autonomy (3.1\%), Misuse of personal data (3.1\%),
    Safety/security/environmental risks (3.0\%).
  \item
    \textbf{2021--2023}: data themes return strongly. Top 5 --- Misuse
    of personal data (6.5\%), Privacy and data security (4.4\%), Lack of
    transparency in decision-making (4.2\%), Insufficient human
    oversight (4.0\%), Mistrust of government data handling (3.8\%).
  \item
    \textbf{2024--2025}: data privacy still top, but joined by themes of
    exclusion and broader oversight. Top 5 --- Misuse of personal data
    (5.4\%), Privacy and data security (4.9\%), Distrust in governance
    and oversight (3.8\%), Exclusion of marginalised groups (3.7\%),
    Lack of transparency in decision-making (3.3\%).
  \end{itemize}
\item
  \textbf{Defensible claim}: AI public discourse has consistently
  centred on data themes, with their relative prominence ebbing and
  flowing as other concerns (governance, equity, exclusion) come into
  focus alongside them.
\item
  \textbf{Caveat}: temporal patterns reflect in part
  what~\emph{kinds}~of dialogues were commissioned in each period. The
  2004--2017 window is dominated by health-data dialogues and other
  early data-governance work; 2018--2020 sees more general-attitudes
  work; 2021--2023 reflects post-pandemic biometrics and trust-related
  dialogues; 2024--2025 is the current AI policy moment. The temporal
  narrative is partly about the commissioning landscape, not just
  shifting public attitudes.
\end{itemize}

\textbf{Figure 4}: heatmap of cluster shares across the four time
windows for the union of clusters appearing in the top 10 of any window.
Rows ordered by which window each cluster peaks in, producing a diagonal
pattern that visualises the content shift.

\includegraphics[width=6.26806in,height=5.66806in]{media/image4.png}

\textbf{Reading the finding}: this is a more constrained temporal claim
than the original v12 draft made, but it is the claim the data supports.
"AI public discourse evolves rapidly" would be an overclaim; "AI public
discourse content mix shifts within a stable spread" is what the
evidence shows.

% ─────────────────────────────────────────────────────────────────────────────
% R4. LIMITATIONS
% Notebook: 05_robustness.ipynb
%   - Inter-rater validation numbers: cell 4 (sample generation); results from run
%   - Prompt sensitivity results: §SECTION 10C (cell 31)
%   - Dendrogram ARI: 02_shared_structure.ipynb cells 4/18
%   - k-sensitivity: 05_robustness.ipynb §SECTION 10/10B
%   - Corpus imbalance and segmentation heterogeneity:
%       00_data_quality.ipynb, 01_processing.ipynb cells 22–24
% ─────────────────────────────────────────────────────────────────────────────
\textbf{R4. Limitations}

\begin{itemize}
\item
  \textbf{Human inter-rater validation}~against a stratified sample of
  50 paragraphs shows 92\% agreement and Cohen's κ = 0.84
  (almost-perfect). The LLM is prone to over-extraction (4 of 50 cases)
  but not under-extraction (0 cases). The bias has a direction: when in
  doubt, the LLM tends to find concerns rather than miss them.
  Validation against a larger sample is a candidate task for the final
  canonical run.
\item
  \textbf{Prompt sensitivity}~analysis confirms direction of headline
  findings holds across V0/V1/V3 variants; magnitudes (specifically the
  Privacy effect) are sensitive to decontextualisation
  strictness.~\emph{(The sensitivity analysis was run against the v16
  corpus; re-running against v19 chunks is a candidate task.)}
\item
  \textbf{Lens-scheme empirical support}: dendrogram-based check
  produces ARI = 0.185 for concerns and 0.150 for benefits, indicating
  the lens scheme is a useful conceptual organisation but does not carve
  at strong joints in the data structure. Some lenses (notably
  Privacy/Consent \& Data) are well-validated by the data; others are
  looser groupings. We present the lens scheme as a deliberate analytic
  choice that organises the data into theoretically meaningful groups,
  with the dendrogram as evidence that some of those groups carve at
  real semantic joints (validating those findings) while others don't
  (noting which lens-level claims to hedge).
\item
  \textbf{Corpus imbalance}: 41 of 66 documents are AI; smaller
  technologies are thinly represented. Inferences about non-AI
  technologies as a whole rest predominantly on the genetic-technology
  and gene-editing subcorpora.
\item
  \textbf{Segmentation heterogeneity}: 15 of 66 documents (mostly recent
  UK AI policy reports produced from HTML originals) lack extractable
  paragraph break structure and were segmented at sentence boundaries
  throughout. A further 43 documents had at least one oversized
  paragraph segmented internally. The remaining 8 documents are pure
  paragraph segmentation. The~chunking\_method~column records this
  per-chunk for audit purposes.
\item
  \textbf{Year coverage gaps}: substantial absences in 2005--2006,
  2011--2014, and 2020 reflect what public dialogues were commissioned
  during those periods rather than any analytical choice.
\item
  Reports were segmented using a hybrid strategy. Where PDF paragraph
  structure was preserved, double-newline paragraph boundaries were
  used. Where paragraphs were unusually long or paragraph boundaries
  were not recoverable, text was split into sentence-based chunks of
  approximately comparable length. Each chunk retained a method flag to
  support quality checks.
\end{itemize}

\textbf{Open items before submission}

(For the author's working list, not for the paper text.)

\begin{enumerate}
\def\labelenumi{\arabic{enumi}.}
\item
  \textbf{Decide on canonical model}: gpt-5-mini (current) vs GPT-5 vs
  both as a robustness check.
\item
  \textbf{Re-run prompt sensitivity against v19 corpus}~if time/budget
  allows; current numbers are from v16 sample.
\item
  \textbf{Expand human validation}~beyond the 50-paragraph proof of
  concept.
\item
  \textbf{Regenerate paper-asset PNGs}~(Figures 1--4) from v19 outputs
  via the in-notebook Paper Assets section.
\end{enumerate}
